% Gerado por ferramentas/gerar.py; não edite à mão.
\section*{Exercícios selecionados}
\addcontentsline{toc}{section}{Exercícios selecionados}
Problemas escolhidos para desenvolver a resolução de problemas. As soluções completas,
com animações dos gráficos, estão no repositório do livro: \url{https://github.com/AleffAzevedo/calculo-1-pt/tree/main/exercicios/cap2}.

\begin{enumerate}[label=\textbf{\thechapter.\arabic*}, leftmargin=2.6em, itemsep=6pt]
\item \textbf{Da secante à tangente.} Calcule $\displaystyle\lim_{h\to0}\frac{\frac{1}{a+h}-\frac1a}{h}$, em que $a\neq0$ é constante.
Interprete o resultado em termos de retas secantes ao gráfico de $f(x)=1/x$.
\item \textbf{O limite de $\sin(ax)/x$.} (a) Use tabelas com $x=\pm0{,}1;\ \pm0{,}01;\ \pm0{,}001$ para estimar
$\lim_{x\to0}\frac{\sin2x}{x}$ e $\lim_{x\to0}\frac{\sin3x}{x}$.
(b) Conjecture o valor de $\lim_{x\to0}\frac{\sin ax}{x}$ para $a>0$ e prove-o com o
Teorema~2.7.
\item \textbf{Conjugado com raiz.} Mostre que a substituição direta leva a uma forma indeterminada e calcule
$\displaystyle\lim_{t\to9}\frac{t-9}{\sqrt t-3}$.
\item \textbf{Fatorando cubos.} Calcule $\displaystyle\lim_{x\to1}\frac{x^3-1}{x^2-1}$.
\item \textbf{Limites laterais infinitos.} Calcule $\displaystyle\lim_{x\to1^-}\frac{2x^2+7x-4}{x^2+x-2}$ e
$\displaystyle\lim_{x\to1^+}\frac{2x^2+7x-4}{x^2+x-2}$.
\item \textbf{Teorema do Confronto.} Calcule $\displaystyle\lim_{\theta\to0}\theta^2\cos\!\left(\frac1\theta\right)$.
\item \textbf{Escolhendo $k$ para a continuidade (I).} Encontre $k$ para que $f(x)=\begin{cases}3x+2, & x<k\\ 2x-3, & k\le x\le8\end{cases}$
seja contínua em todo o seu domínio.
\item \textbf{Escolhendo $k$ para a continuidade (II).} Encontre $k$ para que $f(x)=\begin{cases}e^{kx}, & 0\le x<4\\ x+3, & 4\le x\le8\end{cases}$
seja contínua em $[0,8]$.
\item \textbf{TVI: duas partículas.} Uma partícula se move numa reta com posição $s(t)$ contínua, e $s(2)=5$, $s(5)=2$. Outra
partícula tem posição $h(t)=s(t)-t$. Explique por que existe um instante $t\in[2,5]$ em que
$h(t)=0$. O que isso significa para a primeira partícula?
\item \textbf{TVI: localizando uma raiz.} Use o TVI para decidir se $2^x=x^3$ tem solução em $[1{,}25;\,1{,}375]$ ou em
$[1{,}375;\,1{,}5]$. Justifique para cada intervalo.
\item \textbf{Prova $\varepsilon$-$\delta$.} Use a definição precisa de limite para provar que $\displaystyle\lim_{x\to2}(x^2+2x)=8$.
\item \textbf{Provando que um limite não existe.} Use a definição precisa para provar que $\displaystyle\lim_{x\to1}\frac{|x-1|}{x-1}$ não existe.
\end{enumerate}

\subsection*{Respostas}
{\small\begin{enumerate}[label=\textbf{\thechapter.\arabic*}, leftmargin=2.6em, itemsep=2pt]
\item $-\dfrac{1}{a^2}$: é a inclinação da reta tangente a $y=1/x$ em $x=a$.
\item (a) 2 e 3; (b) $\lim_{x\to0}\frac{\sin ax}{x}=a$.
\item $0/0$; o limite é 6.
\item $\dfrac32$.
\item $-\infty$ pela esquerda e $+\infty$ pela direita.
\item 0, pelo Teorema do Confronto.
\item $k=-5$.
\item $k=\dfrac{\ln7}{4}\approx0{,}4865$.
\item $h(2)=3>0$ e $h(5)=-3<0$; pelo TVI, $h(c)=0$ para algum $c\in[2,5]$, ou seja, $s(c)=c$.
\item Com $g(x)=2^x-x^3$: $g(1{,}25)>0$ e $g(1{,}375)<0$, logo há raiz no 1º intervalo; no 2º, $g$ não muda de sinal e o TVI nada garante.
\item Escolha $\delta=\min\{1,\varepsilon/7\}$.
\item Tome $\varepsilon=\tfrac12$: perto de 1 a função vale $-1$ e $1$, que não cabem juntos em nenhuma faixa de largura 1.
\end{enumerate}}
